\subsection{Collapse variants}
Table~\ref{tab:collapse} varies the collapse window $w$ (exponent free), and fixes $\nu=1$. These are analyses of the saved curves of the main runs, not new training runs. With $\nu=1$ fixed, the mean collapse error of the CNN (\rhval{abl_nu_fixed/cnn/ising/tc_error_fss/mean:3}), MLP (\rhval{abl_nu_fixed/mlp/ising/tc_error_fss/mean:3}) and the Binder reference (\rhval{abl_nu_fixed/binder-chi-reference/ising/tc_error_fss/mean:3}) is lower than with a free exponent in the main table. Paired over seeds (Wilcoxon, fixed minus free) the improvement is detectable for the CNN (mean difference \rhval{analysis/paired-tests/ising/nu1_cnn_meandiff/mean:3}, $p=\rhval{analysis/paired-tests/ising/nu1_cnn_p/mean:3}$, lower on \rhval{analysis/paired-tests/ising/nu1_cnn_wins/mean:0} of ten seeds) and PCA ($p=\rhval{analysis/paired-tests/ising/nu1_pca_p/mean:3}$), but not for the MLP ($p=\rhval{analysis/paired-tests/ising/nu1_mlp_p/mean:2}$) or the Binder reference ($p=\rhval{analysis/paired-tests/ising/nu1_binder_p/mean:2}$). We therefore tie the collapse error to the free exponent only for the CNN; a free exponent running into the search boundary may be the reason, which we did not isolate further. Fixing $\nu=1$ leaves the confusion collapse unchanged. The fixed-$\nu$ collapse error of the CNN is of the same size as its single-size crossing at $L=32$ (\rhval{main/cnn/ising/tc_error_l32/mean:3}); we did not test that difference. We did not test the dependence on the number of sizes. The confusion curves fail in every variant.

\begin{table}[t]\centering\caption{Collapse error $|\hat T_c^{\rm fss}-T_c|$ (mean over seeds) for three windows $w$ with $\nu$ free, with $\nu=1$ fixed ($w=0.6$), and the main-table exponent error.}\label{tab:collapse}
\resizebox{\columnwidth}{!}{\begin{tabular}{lccccc}
\toprule
System & $w{=}0.6$ & $w{=}0.45$ & $w{=}0.3$ & $\nu{=}1$ fixed & $|\nu{-}1|$ free \\
\midrule
CNN & \rhval{main/cnn/ising/tc_error_fss/mean:3} & \rhval{abl_window_0.45/cnn/ising/tc_error_fss/mean:3} & \rhval{abl_window_0.3/cnn/ising/tc_error_fss/mean:3} & \rhval{abl_nu_fixed/cnn/ising/tc_error_fss/mean:3} & \rhval{main/cnn/ising/nu_error_fss/mean:3} \\
MLP & \rhval{main/mlp/ising/tc_error_fss/mean:3} & \rhval{abl_window_0.45/mlp/ising/tc_error_fss/mean:3} & \rhval{abl_window_0.3/mlp/ising/tc_error_fss/mean:3} & \rhval{abl_nu_fixed/mlp/ising/tc_error_fss/mean:3} & \rhval{main/mlp/ising/nu_error_fss/mean:3} \\
Confusion (MLP) & \rhval{main/confusion-mlp/ising/tc_error_fss/mean:3} & \rhval{abl_window_0.45/confusion-mlp/ising/tc_error_fss/mean:3} & \rhval{abl_window_0.3/confusion-mlp/ising/tc_error_fss/mean:3} & \rhval{abl_nu_fixed/confusion-mlp/ising/tc_error_fss/mean:3} & \rhval{main/confusion-mlp/ising/nu_error_fss/mean:3} \\
PCA & \rhval{main/pca/ising/tc_error_fss/mean:3} & \rhval{abl_window_0.45/pca/ising/tc_error_fss/mean:3} & \rhval{abl_window_0.3/pca/ising/tc_error_fss/mean:3} & \rhval{abl_nu_fixed/pca/ising/tc_error_fss/mean:3} & \rhval{main/pca/ising/nu_error_fss/mean:3} \\
Binder/$\chi$ (ref.) & \rhval{main/binder-chi-reference/ising/tc_error_fss/mean:3} & \rhval{abl_window_0.45/binder-chi-reference/ising/tc_error_fss/mean:3} & \rhval{abl_window_0.3/binder-chi-reference/ising/tc_error_fss/mean:3} & \rhval{abl_nu_fixed/binder-chi-reference/ising/tc_error_fss/mean:3} & \rhval{main/binder-chi-reference/ising/nu_error_fss/mean:3} \\
LogReg (Z2-fixed) & \rhval{main/logreg-z2-fixed/ising/tc_error_fss/mean:3} & \rhval{abl_window_0.45/logreg-z2-fixed/ising/tc_error_fss/mean:3} & \rhval{abl_window_0.3/logreg-z2-fixed/ising/tc_error_fss/mean:3} & \rhval{abl_nu_fixed/logreg-z2-fixed/ising/tc_error_fss/mean:3} & \rhval{main/logreg-z2-fixed/ising/nu_error_fss/mean:3} \\
LogReg (raw) & \rhval{main/logreg-raw/ising/tc_error_fss/mean:3} & \rhval{abl_window_0.45/logreg-raw/ising/tc_error_fss/mean:3} & \rhval{abl_window_0.3/logreg-raw/ising/tc_error_fss/mean:3} & \rhval{abl_nu_fixed/logreg-raw/ising/tc_error_fss/mean:3} & \rhval{main/logreg-raw/ising/nu_error_fss/mean:3} \\
\bottomrule
\end{tabular}
}
\end{table}

\subsection{H6: sensitivity to the training window and to the training-set size}
Table~\ref{tab:sweeps} (five seeds, 0--4, default rows from the main group for the same seeds) varies the half-width $\delta$ of the excluded window and the number of samples per chain used for training (3 chains $\times$ 10, 25 or 50 samples per temperature). In the first training-size sweep the number of epochs is fixed, so fewer samples also means proportionally fewer gradient updates; we therefore added a second sweep (\texttt{sweep\_nsamp\_ep}, same five seeds) in which the number of epochs is scaled so that samples $\times$ epochs equals that of the default.
\emph{Window.} Between $\delta=0.15$ and $0.6$ the $L=32$ error changes only moderately (ratio column). The largest ratios are \rhval{analysis/sweep-summary/ising/mlp_margin0p6_ratio_l32/mean:2} (MLP) and \rhval{analysis/sweep-summary/ising/cnn_margin0p6_ratio_l32/mean:2} (CNN) at $\delta=0.6$, i.e.\ the windows can be shrunk or enlarged by a factor two with an effect below the registered factor-two threshold.
\emph{Training-set size.} Using only $10$ samples per chain multiplies the $L=32$ error by \rhval{analysis/sweep-summary/ising/mlp_nsamp10_ratio_l32/mean:2} (MLP) and \rhval{analysis/sweep-summary/ising/cnn_nsamp10_ratio_l32/mean:2} (CNN), above the registered threshold of two at a fixed number of epochs. With the number of gradient updates held fixed (10 samples per chain, \rhval{const/nsamp_ep10_epochs} epochs) the ratio is \rhval{analysis/sweep-summary/ising/mlp_nsamp_ep10_ratio_l32/mean:2} for the MLP, still above two, but \rhval{analysis/sweep-summary/ising/cnn_nsamp_ep10_ratio_l32/mean:2} for the CNN, whose apparent sensitivity (above) is thus an optimisation-budget effect. \textbf{H6 is therefore refuted for the MLP and supported for the CNN} as far as training-set size is concerned (the registered statement is about both; with five seeds and no further test, the verdict for the MLP rests on a ratio of means). The collapse error is much less affected by the training size, and for the MLP at $\delta=0.15$ it is smaller than the default (ratio \rhval{analysis/sweep-summary/ising/mlp_margin0p15_ratio_fss/mean:2}); with five seeds and a std comparable to the means, we do not interpret these differences. Fig.~\ref{fig:sweeps} shows the single-size error at $L=32$.

\begin{table}[t]\centering\caption{Sensitivity (five seeds, mean): $|\hat T_c-T_c|$ at $L=16$, $L=32$, after collapse, and the ratio of the $L=32$ error to the default setting (margin $0.3$, $50$ samples/chain). ``scaled ep.'': epochs scaled to keep the number of gradient updates fixed.}\label{tab:sweeps}
\resizebox{\columnwidth}{!}{\begin{tabular}{llcccc}
\toprule
Model & setting & $L{=}16$ & $L{=}32$ & collapse & ratio$_{32}$ \\
\midrule
CNN & margin 0.15 & \rhval{analysis/sweep-summary/ising/cnn_margin0p15_l16_mean/mean:3} & \rhval{analysis/sweep-summary/ising/cnn_margin0p15_l32_mean/mean:3} & \rhval{analysis/sweep-summary/ising/cnn_margin0p15_fss_mean/mean:3} & \rhval{analysis/sweep-summary/ising/cnn_margin0p15_ratio_l32/mean:2} \\
CNN & margin 0.3 & \rhval{analysis/sweep-summary/ising/cnn_margin0p3_l16_mean/mean:3} & \rhval{analysis/sweep-summary/ising/cnn_margin0p3_l32_mean/mean:3} & \rhval{analysis/sweep-summary/ising/cnn_margin0p3_fss_mean/mean:3} & \rhval{analysis/sweep-summary/ising/cnn_margin0p3_ratio_l32/mean:2} \\
CNN & margin 0.45 & \rhval{analysis/sweep-summary/ising/cnn_margin0p45_l16_mean/mean:3} & \rhval{analysis/sweep-summary/ising/cnn_margin0p45_l32_mean/mean:3} & \rhval{analysis/sweep-summary/ising/cnn_margin0p45_fss_mean/mean:3} & \rhval{analysis/sweep-summary/ising/cnn_margin0p45_ratio_l32/mean:2} \\
CNN & margin 0.6 & \rhval{analysis/sweep-summary/ising/cnn_margin0p6_l16_mean/mean:3} & \rhval{analysis/sweep-summary/ising/cnn_margin0p6_l32_mean/mean:3} & \rhval{analysis/sweep-summary/ising/cnn_margin0p6_fss_mean/mean:3} & \rhval{analysis/sweep-summary/ising/cnn_margin0p6_ratio_l32/mean:2} \\
CNN & samples/chain 10 & \rhval{analysis/sweep-summary/ising/cnn_nsamp10_l16_mean/mean:3} & \rhval{analysis/sweep-summary/ising/cnn_nsamp10_l32_mean/mean:3} & \rhval{analysis/sweep-summary/ising/cnn_nsamp10_fss_mean/mean:3} & \rhval{analysis/sweep-summary/ising/cnn_nsamp10_ratio_l32/mean:2} \\
CNN & samples/chain 25 & \rhval{analysis/sweep-summary/ising/cnn_nsamp25_l16_mean/mean:3} & \rhval{analysis/sweep-summary/ising/cnn_nsamp25_l32_mean/mean:3} & \rhval{analysis/sweep-summary/ising/cnn_nsamp25_fss_mean/mean:3} & \rhval{analysis/sweep-summary/ising/cnn_nsamp25_ratio_l32/mean:2} \\
CNN & samples/chain 50 & \rhval{analysis/sweep-summary/ising/cnn_nsamp50_l16_mean/mean:3} & \rhval{analysis/sweep-summary/ising/cnn_nsamp50_l32_mean/mean:3} & \rhval{analysis/sweep-summary/ising/cnn_nsamp50_fss_mean/mean:3} & \rhval{analysis/sweep-summary/ising/cnn_nsamp50_ratio_l32/mean:2} \\
CNN & samples, scaled ep. 10 & \rhval{analysis/sweep-summary/ising/cnn_nsamp_ep10_l16_mean/mean:3} & \rhval{analysis/sweep-summary/ising/cnn_nsamp_ep10_l32_mean/mean:3} & \rhval{analysis/sweep-summary/ising/cnn_nsamp_ep10_fss_mean/mean:3} & \rhval{analysis/sweep-summary/ising/cnn_nsamp_ep10_ratio_l32/mean:2} \\
CNN & samples, scaled ep. 25 & \rhval{analysis/sweep-summary/ising/cnn_nsamp_ep25_l16_mean/mean:3} & \rhval{analysis/sweep-summary/ising/cnn_nsamp_ep25_l32_mean/mean:3} & \rhval{analysis/sweep-summary/ising/cnn_nsamp_ep25_fss_mean/mean:3} & \rhval{analysis/sweep-summary/ising/cnn_nsamp_ep25_ratio_l32/mean:2} \\
\midrule
MLP & margin 0.15 & \rhval{analysis/sweep-summary/ising/mlp_margin0p15_l16_mean/mean:3} & \rhval{analysis/sweep-summary/ising/mlp_margin0p15_l32_mean/mean:3} & \rhval{analysis/sweep-summary/ising/mlp_margin0p15_fss_mean/mean:3} & \rhval{analysis/sweep-summary/ising/mlp_margin0p15_ratio_l32/mean:2} \\
MLP & margin 0.3 & \rhval{analysis/sweep-summary/ising/mlp_margin0p3_l16_mean/mean:3} & \rhval{analysis/sweep-summary/ising/mlp_margin0p3_l32_mean/mean:3} & \rhval{analysis/sweep-summary/ising/mlp_margin0p3_fss_mean/mean:3} & \rhval{analysis/sweep-summary/ising/mlp_margin0p3_ratio_l32/mean:2} \\
MLP & margin 0.45 & \rhval{analysis/sweep-summary/ising/mlp_margin0p45_l16_mean/mean:3} & \rhval{analysis/sweep-summary/ising/mlp_margin0p45_l32_mean/mean:3} & \rhval{analysis/sweep-summary/ising/mlp_margin0p45_fss_mean/mean:3} & \rhval{analysis/sweep-summary/ising/mlp_margin0p45_ratio_l32/mean:2} \\
MLP & margin 0.6 & \rhval{analysis/sweep-summary/ising/mlp_margin0p6_l16_mean/mean:3} & \rhval{analysis/sweep-summary/ising/mlp_margin0p6_l32_mean/mean:3} & \rhval{analysis/sweep-summary/ising/mlp_margin0p6_fss_mean/mean:3} & \rhval{analysis/sweep-summary/ising/mlp_margin0p6_ratio_l32/mean:2} \\
MLP & samples/chain 10 & \rhval{analysis/sweep-summary/ising/mlp_nsamp10_l16_mean/mean:3} & \rhval{analysis/sweep-summary/ising/mlp_nsamp10_l32_mean/mean:3} & \rhval{analysis/sweep-summary/ising/mlp_nsamp10_fss_mean/mean:3} & \rhval{analysis/sweep-summary/ising/mlp_nsamp10_ratio_l32/mean:2} \\
MLP & samples/chain 25 & \rhval{analysis/sweep-summary/ising/mlp_nsamp25_l16_mean/mean:3} & \rhval{analysis/sweep-summary/ising/mlp_nsamp25_l32_mean/mean:3} & \rhval{analysis/sweep-summary/ising/mlp_nsamp25_fss_mean/mean:3} & \rhval{analysis/sweep-summary/ising/mlp_nsamp25_ratio_l32/mean:2} \\
MLP & samples/chain 50 & \rhval{analysis/sweep-summary/ising/mlp_nsamp50_l16_mean/mean:3} & \rhval{analysis/sweep-summary/ising/mlp_nsamp50_l32_mean/mean:3} & \rhval{analysis/sweep-summary/ising/mlp_nsamp50_fss_mean/mean:3} & \rhval{analysis/sweep-summary/ising/mlp_nsamp50_ratio_l32/mean:2} \\
MLP & samples, scaled ep. 10 & \rhval{analysis/sweep-summary/ising/mlp_nsamp_ep10_l16_mean/mean:3} & \rhval{analysis/sweep-summary/ising/mlp_nsamp_ep10_l32_mean/mean:3} & \rhval{analysis/sweep-summary/ising/mlp_nsamp_ep10_fss_mean/mean:3} & \rhval{analysis/sweep-summary/ising/mlp_nsamp_ep10_ratio_l32/mean:2} \\
MLP & samples, scaled ep. 25 & \rhval{analysis/sweep-summary/ising/mlp_nsamp_ep25_l16_mean/mean:3} & \rhval{analysis/sweep-summary/ising/mlp_nsamp_ep25_l32_mean/mean:3} & \rhval{analysis/sweep-summary/ising/mlp_nsamp_ep25_fss_mean/mean:3} & \rhval{analysis/sweep-summary/ising/mlp_nsamp_ep25_ratio_l32/mean:2} \\
\bottomrule
\end{tabular}
}
\end{table}

\begin{figure}[t]\centering\includegraphics[width=\columnwidth]{fig_sweeps.pdf}
\caption{$L=32$ error versus window half-width (left) and samples per chain (right); mean $\pm$ std over seeds 0--4. Dashed (right): epochs scaled to a constant number of gradient updates.}\label{fig:sweeps}\end{figure}
